\documentclass{article}
\usepackage{mathnotes}

\title{Elementary Cellular Automata}
\description{One-dimensional cellular automata with binary states and local rules, exploring Wolfram's 256 elementary rules including chaotic patterns, fractals, and Turing completeness with interactive demonstrations.}

\begin{document}

\section{Elementary Cellular Automata}

Elementary cellular automata are one-dimensional cellular automata with two possible states (0 or 1) and rules that depend on the cell's current state and the states of its two immediate neighbors. These simple systems were extensively studied by Stephen Wolfram and can produce surprisingly complex behavior.

\subsection{How They Work}

Each cell's next state is determined by:

\begin{itemize}
\item Its current state
\item The state of its left neighbor
\item The state of its right neighbor
\end{itemize}

Since each of these can be either 0 or 1, there are $2^3 = 8$ possible neighborhood configurations. A rule specifies the next state for each configuration, giving us $2^8 = 256$ possible rules.

\subsection{Rule Numbering}

Rules are numbered from 0 to 255 based on their binary representation. For example, Rule 30:

\begin{tabular}{lllllllll}
Pattern & 111 & 110 & 101 & 100 & 011 & 010 & 001 & 000 \\
\hline
Result & 0 & 0 & 0 & 1 & 1 & 1 & 1 & 0 \\
\end{tabular}

The binary number 00011110 equals 30 in decimal, hence "Rule 30".

\subsection{Interactive Demonstration}

The demonstration below allows you to explore different elementary cellular automata rules. You can:

\begin{itemize}
\item Select different rules using the dropdown or enter a custom rule number
\item Adjust the simulation speed
\item See how different initial conditions evolve
\end{itemize}

\includedemo{elementary-cellular-automata}

\subsection{Notable Rules}

\subsubsection{Rule 30}
Produces chaotic, seemingly random patterns despite its simple, deterministic nature. Used in Mathematica's random number generator.

\subsubsection{Rule 90}
Creates the Sierpiński triangle fractal pattern. Equivalent to Pascal's triangle modulo 2.

\subsubsection{Rule 110}
Proven to be Turing complete, meaning it can perform any computation given the right initial conditions.

\subsubsection{Rule 184}
Models traffic flow and ballistic annihilation. Useful for studying particle systems.

\subsection{Mathematical Properties}

Elementary cellular automata exhibit various interesting properties:

\begin{enumerate}
\item \textbf{Reversibility}: Some rules (like Rule 51) are reversible - you can uniquely determine previous states from current ones
\item \textbf{Conservation}: Some rules conserve the number of live cells
\item \textbf{Symmetry}: Many rules produce symmetric patterns when started from symmetric initial conditions
\item \textbf{Fractals}: Several rules generate self-similar fractal patterns
\end{enumerate}

\subsection{Applications}

Despite their simplicity, elementary cellular automata have applications in:

\begin{itemize}
\item Random number generation
\item Cryptography
\item Modeling physical phenomena
\item Pattern generation in computer graphics
\item Understanding emergence and complexity
\end{itemize}

\subsection{Further Reading}

\begin{itemize}
\item Wolfram, S. (2002). \emph{A New Kind of Science}
\item Cook, M. (2004). "Universality in Elementary Cellular Automata"
\item \href{https://mathworld.wolfram.com/ElementaryCellularAutomaton.html}{Wolfram MathWorld: Elementary Cellular Automaton}
\end{itemize}

\end{document}
